\section{Conclusion}
In this work, we proposed the Adaptive Retrieval-Augmented Generation framework, referred to as Adaptive-RAG, to handle queries of various complexities. Specifically, Adaptive-RAG is designed to dynamically adjust its query handling strategies in the unified retrieval-augmented LLM based on the complexity of queries that they encounter, which spans across a spectrum of the non-retrieval-based approach for the most straightforward queries, to the single-step approach for the queries of moderate complexity, and finally to the multi-step approach for the complex queries. The core step of our Adaptive-RAG lies in determining the complexity of the given query, which is instrumental in selecting the most suitable strategy for its answer. To operationalize this process, we trained a smaller Language Model with query-complexity pairs, which are automatically annotated from the predicted outcomes and the inductive biases in datasets. We validated our Adaptive-RAG on a collection of open-domain QA datasets, covering the multiple query complexities including both the single- and multi-hop questions. The results demonstrate that our Adaptive-RAG enhances the overall accuracy and efficiency of QA systems, allocating more resources to handle complex queries while efficiently handling simpler queries, compared to the existing one-size-fits-all approaches that tend to be either minimalist or maximalist over varying query complexities.

